\section[Cyclotomic double relations: exact identities and formal obstructions]{Cyclotomic double relations:\\exact identities and formal obstructions}
\label{sec:gaussian}

Set
\[
 g_{a,b}=\Impart\Li_{a,b}(\ii,1),\quad
 A_{a,b}=\Impart\Li_{a,b}(\ii,\ii),\quad
 K_{a,b}=\Impart\Li_{a,b}(\ii,-1).
\]
The manuscript's mixed harmonic sums are
\[
 S_p=\sum_{n\geq0}\frac{(-1)^nH_n}{(2n+1)^p},\qquad p\geq1.
\]
For \(p>1\) this converges absolutely; at \(p=1\) it converges
conditionally. The all-odd-index evaluation of \(S_p\) has already
been proved in the repository's continuations. Our questions concern
even indices and what exact relations can establish about them.

\subsection{Two exact weight-five reductions}

\begin{proposition}[Two coordinates suffice for the weight-five \(g\)-family]
\label{prop:two-g}
The four values of weight five satisfy
\begin{align}
 g_{2,3}+3g_{3,2}+6g_{4,1}
   &=-\frac3{32}G\zeta(3)-\frac{\pi^5}{1536},
     \label{eq:g-five-first}\\
 g_{1,4}+g_{2,3}+g_{3,2}+2g_{4,1}
   &=-\frac12\beta(4)\log2-\frac{7\pi^5}{46080}.
     \label{eq:g-five-second}
\end{align}
Thus they lie in the span of any two suitably chosen coordinates,
for example \(g_{4,1},g_{2,3}\), together with the displayed products
of single values. This is a spanning statement, not an independence
statement.
\end{proposition}

\begin{proof}
Shuffle the two iterated-integral words for
\(\Li_2(\ii)\Li_3(\ii)\). Their imaginary part is
\(g_{2,3}+3g_{3,2}+6g_{4,1}\).
Since
\[
 \Li_2(\ii)=-\frac{\zeta(2)}8+\ii G,\qquad
 \Li_3(\ii)=-\frac{3\zeta(3)}{32}+\ii\frac{\pi^3}{32},
\]
its product gives \eqref{eq:g-five-first}.
Likewise, the shuffle of \(\Li_1(\ii)\Li_4(\ii)\) has imaginary part
\(g_{1,4}+g_{2,3}+g_{3,2}+2g_{4,1}\). Use
\[
 \Li_1(\ii)=-\tfrac12\log2+\ii\tfrac\pi4,\qquad
 \Li_4(\ii)=-\tfrac7{128}\zeta(4)+\ii\beta(4).
\]
The two rows have rank two already in their \(g\)-coefficients.
\end{proof}

The manuscript's displayed sporadic row
\[
 576g_{4,1}+288g_{3,2}+736g_{2,3}+960g_{1,4}
 =21G\zeta(3)-480\beta(4)\log2
\]
is \(-224\) times \eqref{eq:g-five-first} plus \(960\) times
\eqref{eq:g-five-second}. It therefore has a direct proof from
these two elementary shuffles. The manuscript's \emph{at most three}
statement can be strengthened to \emph{at most two}. Existing
continuations already address proofs of the sporadic relation; the
point here is the sharper two-row organization.

\subsection{An all-exponent mixed endpoint identity}

\begin{theorem}[Gaussian mixed endpoint]\label{thm:mixed-endpoint}
For every integer \(p\geq1\),
\begin{equation}\label{eq:mixed-endpoint}
 K_{p,1}+(-1)^{p-1}A_{1,p}
 =\beta(p+1)-2\log2\,\beta(p)
       +\sum_{j=2}^{p}(-1)^j\eta(j)\beta(p+1-j).
\end{equation}
Moreover,
\begin{equation}\label{eq:harmonic-split}
 S_p=g_{p,1}+K_{p,1}.
\end{equation}
\end{theorem}

\begin{proof}
All the boundary values used here converge. First-index-one values
with first argument \(\ii\) can be obtained by Dirichlet's test or
radial Abel limits. The shuffle and stuffle identities may therefore
be justified by convergent iterated integrals and Abel limiting.

Fix \(w=p+1\), and define homogeneous generating polynomials of
degree \(w-2\):
\[
 \mathscr A(X,Y)=\sum_{a+b=w}A_{a,b}X^{a-1}Y^{b-1},
 \quad
 \mathscr K(X,Y)=\sum_{a+b=w}K_{a,b}X^{a-1}Y^{b-1}.
\]
Let \(\mathscr V\) denote the same polynomial for
\(\Impart\Li_{a,b}(-1,\ii)\), and put
\[
 P(X,Y)=\sum_{a+b=w}
  \Impart[\Li_a(-1)\Li_b(\ii)]X^{a-1}Y^{b-1},
 \qquad
 H(X,Y)=\sum_{a+b=w}X^{a-1}Y^{b-1}.
\]
Stuffle and shuffle respectively give
\[
 \mathscr V(X,Y)+\mathscr K(Y,X)=P(X,Y)+\beta(w)H(X,Y),
\]
\[
 \mathscr A(X+Y,X)-\mathscr V(X+Y,Y)=P(X,Y).
\]
The collision term in the first equation is
\(-\Impart\Li_w(-\ii)=\beta(w)\).
Eliminating \(\mathscr V\) yields
\[
 \mathscr K(U,V)=P(V-U,U)+P(V,U)+\beta(w)H(V,U)
                  -\mathscr A(V,V-U).
\]
Set \(V=0\) and compare coefficients of \(U^{w-2}\).
The evaluation \(\Li_j(-1)=-\eta(j)\) gives
\eqref{eq:mixed-endpoint}, including the two occurrences of
\(-\log2\,\beta(p)\).

For \eqref{eq:harmonic-split}, taking imaginary parts of the outer
index in \(\Li_{p,1}(\ii,1)+\Li_{p,1}(\ii,-1)\) retains indices
\(2n+1\), with sign \((-1)^n\). Their inner sum is
\[
 \sum_{m=1}^{2n}\frac{1+(-1)^m}{m}=H_n.
\]
Absolute convergence proves the result directly for \(p>1\).
An Abel limit gives the endpoint \(p=1\).
\end{proof}

At \(p=4\) this specializes to
\begin{equation}\label{eq:K41}
 K_{4,1}=A_{1,4}+\frac{191\pi^5}{23040}
                  -\frac34G\zeta(3)-2\beta(4)\log2.
\end{equation}

\subsection{Inverse-argument mixed values and a source correction}

A useful regularization identity also removes a specific proposed
mixed survivor from the manuscript. Inverse-argument double reductions
are part of the existing repository programme; the following short
derivation makes the correction explicit.

\begin{proposition}[Inverse-argument endpoint]\label{prop:inverse-mixed}
For \(p\geq1\) and \(z=\ii\) or \(e^{2\pi\ii/3}\),
\begin{align}
 \Impart\Li_{p,1}(z,z^{-1})
 &=\sum_{a=1}^{p}\Impart\Li_{a,p+1-a}(z,1)
     +p\,\Impart\Li_{p+1}(z)\notag\\
 &\quad-\sum_{j=2}^{p}\zeta(j)\Impart\Li_{p+1-j}(z).
 \label{eq:inverse-mixed}
\end{align}
\end{proposition}

\begin{proof}
Give \(\Li_1(1)\) its usual real regularization parameter \(T\).
The shuffle expansion of \(T\Li_p(z)\) is
\[
 \Li_{p,1}(z,z^{-1})+\sum_{a=1}^{p}\Li_{a,p+1-a}(1,z),
\]
and its stuffle expansion is
\[
 \Li_{1,p}(1,z)+\Li_{p,1}(z,1)+\Li_{p+1}(z).
\]
The divergent leading term cancels on subtraction. For \(2\leq a\leq p\)
use the convergent stuffle
\[
 \Li_{a,p+1-a}(1,z)
 =\zeta(a)\Li_{p+1-a}(z)
       -\Li_{p+1-a,a}(z,1)-\Li_{p+1}(z).
\]
Taking imaginary parts gives \eqref{eq:inverse-mixed}.
For \(p=1\) there is only one divergent factor because \(z\ne1\);
the same cancellation proves the formula.
\end{proof}

Substitution of the already proved Gaussian weight-six \(g\)-values
gives the particularly short identity
\begin{equation}\label{eq:mixed-weight-six}
 \boxed{\displaystyle
 \Impart\Li_{5,1}(\ii,-\ii)
 =3\beta(6)-\frac{\pi^2}{6}\beta(4)
            -\frac{\pi^4}{90}G-\frac{5\pi^5}{3072}\log2.}
\end{equation}
Thus this imaginary weight-six value is in the depth-one product
basket. The candidate-survivor wording in Chapter~4 should be removed.

For \(\omega=e^{2\pi\ii/3}\), define
\(C_s=\Impart\Li_s(\omega)\); this is the sine-part convention for every
integer \(s\). Then
\begin{align}
 \Impart\Li_{5,1}(\omega,\omega^2)
 &=\sum_{a=1}^{5}\Impart\Li_{a,6-a}(\omega,1)
     +5C_6-\frac{\pi^2}{6}C_4-\frac{\pi^4}{90}C_2\notag\\
 &\quad-\frac{2\pi^3}{81}\zeta(3)-\frac{\pi}{6}\zeta(5).
 \label{eq:eisenstein-weight-six}
\end{align}
Here \(C_1=\pi/6\) and \(C_3=2\pi^3/81\).
It follows that the corresponding level-three imaginary value
contributes no new direction modulo the stated \(g\)-family and
single products. This conclusion says nothing about the manuscript's
real weight-five candidate.

\subsection{A precisely defined formal presentation}

Fix a level \(L=3\) or \(4\) and a weight \(w\geq2\). Let
\(\omega=e^{2\pi\ii/L}\).
Introduce rational formal symbols
\[
 D_a(r,s)\quad(1\leq a<w,\quad r,s\in\Z/L\Z)
\]
for imaginary parts of \(\Li_{a,w-a}(\omega^r,\omega^s)\).
Impose conjugation \(D_a(-r,-s)=-D_a(r,s)\).
Symbols fixed by this involution vanish. Put all single values and
imaginary parts of products of two single values in the background,
and also set \(D_a(1,0)=0\). The latter is precisely quotienting by
the chosen \(g\)-family.

The only further rows imposed are, for \(p+q=w\),
\begin{align}
 D_p(r,s)+D_q(s,r)&=0,\label{eq:formal-stuffle}\\
 \sum_{j=0}^{p-1}\binom{q+j-1}{j}D_{q+j}(s,r-s)
 +\sum_{j=0}^{q-1}\binom{p+j-1}{j}D_{p+j}(r,s-r)&=0.
 \label{eq:formal-shuffle}
\end{align}
They are the stuffle and shuffle rows of two single values, with their
one-leading-one regularizations. A regularized product involving the
real parameter \(T\) belongs to the background. At \(w=2\), the only
two-divergent-factor correction is real and disappears on taking
imaginary parts. Call the resulting quotient \(V_{L,w}\).

This is \emph{not} the full cyclotomic double-shuffle algebra. Products
involving a depth-two factor, relations lifted through higher depths,
additional functional equations, and motivic period assertions have
not been imposed. In particular, Zhao's standard relations
\cite{Zhao} are substantially broader than this presentation.
The theorem below computes \(V_{L,w}\), not a dimension of actual
numbers.

For a scalar assignment to the symbols, introduce
\[
 \mathscr F^{r,s}(X,Y)=
 \sum_{a=1}^{w-1}D_a(r,s)X^{a-1}Y^{w-a-1}.
\]
The conditions that the assignment annihilate all rows are exactly
\begin{align}
 \mathscr F^{r,s}(X,Y)+\mathscr F^{s,r}(Y,X)&=0,
 \label{eq:poly-stuffle}\\
 \mathscr F^{s,r-s}(X+Y,X)+
       \mathscr F^{r,s-r}(X+Y,Y)&=0.
 \label{eq:poly-shuffle}
\end{align}
These follow by comparing binomial coefficients after substitution.
They describe the dual of \(V_{L,w}\), so their solution dimension is
the required quotient dimension.

\subsection{All-weight dimension theorems}

\begin{theorem}[Level four]\label{thm:level-four}
The formal quotient satisfies
\[
 \dim_\Q V_{4,w}=
 \begin{cases}
 0,&w\ \text{even},\\
 (w-1)/2,&w\ \text{odd}.
 \end{cases}
 \qquad
 \sum_{w\geq2}\dim V_{4,w}\,t^w=\frac{t^3}{(1-t^2)^2}.
\]
At odd weight, its dual is parameterized by arbitrary antisymmetric
homogeneous polynomials of degree \(w-2\).
\end{theorem}

\begin{proof}
Up to conjugation, the six color pairs are
\((0,1),(1,0),(1,1),(1,2),(1,3),(2,1)\).
Write their polynomials as \(U,G,A,K,M,V\), respectively, and let
\(n=w-2\). The background condition gives \(G=0\);
stuffle gives
\[
 U=0,\quad A(Y,X)=-A(X,Y),\quad
 V(X,Y)=-K(Y,X),\quad M(Y,X)=M(X,Y).
\]
The shuffle with colors \((1,0)\) forces \(M=0\).
The shuffle with colors \((2,1)\) gives
\[
 A(X+Y,X)+K(Y,X+Y)=0,
\]
hence \(K(X,Y)=-A(Y,Y-X)\).
The shuffle with colors \((1,3)\) is
\(K(X,Y)=K(X,X-Y)\).
After substitution and antisymmetry this condition is
\(A=(-1)^{n+1}A\).

For even \(n\) all polynomials therefore vanish.
For odd \(n\), any antisymmetric \(A\), with \(K,V\) as above and
\(U=G=M=0\), satisfies the rows. To check exhaustiveness, the
unordered color classes up to conjugation are
\[
 (0,0),(0,1),(0,2),(1,1),(1,2),(1,3),(2,2).
\]
The real classes vanish, and the remaining classes are precisely the
conditions just used, together with their transposes. Thus there are
no unexamined rows.
A basis for \(A\) at odd \(w\) is
\[
 X^{a-1}Y^{w-a-1}-X^{w-a-1}Y^{a-1},
 \qquad 1\leq a\leq(w-1)/2.
\]
Counting these basis vectors and summing over odd weights proves
the formulas.
\end{proof}

\begin{theorem}[Level three]\label{thm:level-three}
The formal quotient satisfies
\[
 \dim_\Q V_{3,w}=
 \begin{cases}
 0,&w\ \text{even},\\
 \lfloor(w+1)/6\rfloor,&w\ \text{odd},
 \end{cases}
 \qquad
 \sum_{w\geq2}\dim V_{3,w}\,t^w
       =\frac{t^5}{(1-t^2)(1-t^6)}.
\]
Writing \(\Delta=XY(X-Y)\), \(Q=X^2-XY+Y^2\), the dual is parameterized
by the homogeneous polynomials
\[
 A(X,Y)=\Delta\,P(Q,\Delta^2)
\]
of degree \(w-2\).
\end{theorem}

\begin{proof}
Up to conjugation the four pairs are
\((0,1),(1,0),(1,1),(1,2)\).
As before the first two vanish and the shuffle with a zeta factor kills
the inverse-argument polynomial. The remaining polynomial \(A\)
satisfies exactly
\[
 A(Y,X)=-A(X,Y),\qquad A(X,X-Y)=A(X,Y).
\]
Antisymmetry implies divisibility by \(X-Y\); the second symmetry then
forces divisibility by \(Y\), and antisymmetry forces divisibility by
\(X\). Hence \(A=\Delta B\). The polynomial \(\Delta\) has the same
two transformation signs as \(A\), so \(B\) is invariant under swapping
\(X,Y\) and under \((X,Y)\mapsto(X,X-Y)\).

On the plane \(a+b+c=0\), with
\((a,b,c)=(X,-Y,Y-X)\), the second transformation swaps \(b,c\);
the first is global negation followed by the transposition of \(a,b\).
The generated group contains global negation, obtained by cubing
the product, and all permutations. The elementary symmetric-polynomial
theorem therefore identifies its invariant ring with
\[
 \Q[e_2,e_3^2]=\Q[Q,\Delta^2].
\]
Conversely every \(\Delta P(Q,\Delta^2)\) has the required symmetries,
and direct substitution verifies all the original color rows.
Counting monomials of degrees \(3+2a+6b=w-2\) gives the result.
Explicitly, at odd \(w\geq5\) a basis is
\[
 \Delta^{2b+1}Q^{(w-5)/2-3b},
 \qquad 0\leq b\leq\lfloor(w-5)/6\rfloor .
\]
\end{proof}

For comparison, the exact dimensions at the first few odd weights are
\[
\begin{array}{c|rrrrrrr}
 w&3&5&7&9&11&13&15\\ \hline
 \dim V_{3,w}&0&1&1&1&2&2&2\\
 \dim V_{4,w}&1&2&3&4&5&6&7
\end{array}
\]
and all even-weight dimensions are zero. The exact matrix script
checks both formulas through weight 16, including that the displayed
polynomial bases annihilate every raw row and have the stated rank.

\subsection{A separating functional for the \texorpdfstring{$S_4$}{S4} conjecture}

At every odd weight \(w\geq3\), choose
\(A(X,Y)=(Y-X)^{w-2}\) in the level-four theorem.
Then \(K(X,Y)=X^{w-2}\), \(V(X,Y)=-Y^{w-2}\).
This gives a rational linear functional \(\ell\) that annihilates every
defining row and every background value:
\begin{align}
 \ell(D_a(1,1))&=(-1)^{a+1}\binom{w-2}{a-1},\notag\\
 \ell(D_{w-1}(1,2))&=1,\qquad
 \ell(D_1(2,1))=-1. \label{eq:separating}
\end{align}
All other entries vanish, apart from conjugates with opposite signs.

At weight five its nonzero entries are particularly transparent:
\[
\begin{array}{c|rrrrrr}
 \text{symbol}&A_{1,4}&A_{2,3}&A_{3,2}&A_{4,1}
       &\Impart\Li_{1,4}(-1,\ii)&K_{4,1}\\ \hline
 \ell&1&-3&3&-1&-1&1
\end{array}
\]
The polynomial proof verifies all weights; the delivered certificate
also checks this six-entry vector against every exact weight-five row.

\begin{conjecture}[The manuscript's \(S_4\) identity]\label{conj:S4}
The following identity remains unproved in this continuation:
\begin{align}
 S_4&=\frac{4g_{4,1}-3g_{3,2}-9g_{2,3}}7
       +\frac{\pi^5}{224}-\frac{27}{224}G\zeta(3)
       -2\beta(4)\log2. \label{eq:S4-original}
\end{align}
\end{conjecture}

The exact identity \eqref{eq:K41} makes this equivalent to
\begin{equation}\label{eq:S4-A}
 161280A_{1,4}+69120(g_{4,1}+g_{3,2}+3g_{2,3})
       +617\pi^5-101520G\zeta(3)=0.
\end{equation}
The functional \(\ell\) sends its left side to \(161280\).
Consequently this relation does not follow from
\eqref{eq:formal-stuffle}--\eqref{eq:formal-shuffle} and the stated
background alone. This is an exact obstruction to a specific proof
strategy. Additional true analytic relations may still prove it.

Using \eqref{eq:g-five-first}, an equivalent version with only two
\(g\)-coordinates is
\begin{equation}\label{eq:S4-two}
 S_4=\frac{10g_{4,1}-8g_{2,3}}7
       +\frac{7\pi^5}{1536}-\frac3{28}G\zeta(3)
       -2\beta(4)\log2.
\end{equation}
All three displayed candidates are equivalent by proved identities;
none is silently promoted to a theorem.

\subsection{A new weight-seven candidate}

\begin{conjecture}[A reduction for \(S_6\)]\label{conj:S6}
With the same definitions,
\begin{align}
 S_6={}&\frac{722}{527}g_{6,1}
       +\frac{40}{527}g_{4,3}
       -\frac{128}{155}g_{2,5}
       +\frac{15191}{28569600}\pi^7\notag\\
      &-\frac3{124}G\zeta(5)
       -\frac{2373}{10540}\beta(4)\zeta(3)
       -2\beta(6)\log2. \label{eq:S6}
\end{align}
\end{conjecture}

The candidate was found in the ordered basket
\[
 (S_6,g_{6,1},g_{4,3},g_{2,5},\pi^7,
          G\zeta(5),\beta(4)\zeta(3),\beta(6)\log2).
\]
Its primitive integer relation vector is
\[
 \begin{split}
 (&485683200,-665395200,-36864000,401080320,\\
  &-258247,11750400,109347840,971366400).
 \end{split}
\]
A 100-digit search used coefficient bound \(10^{12}\) and tolerance
\(10^{-90}\). After the coefficients were frozen, reevaluation by
Mellin integrals at 220 working digits gave an integer-normalized
residual of approximately \(-2.315\times10^{-213}\).
An independent Euler transformation of the defining alternating
series, using 900 terms and 300 working digits, gave approximately
\(2.854\times10^{-263}\). The script records the full values and
the exact integer vector. These are numerical evidence, not a proof
or an interval certificate.

For the second method, the defining series are
\[
 g_{a,b}=\sum_{n\geq0}
       \frac{(-1)^nH_{2n}^{(b)}}{(2n+1)^a},
 \qquad H_m^{(b)}=\sum_{j=1}^{m}j^{-b},
\]
and the analogous series with \(H_n\) for \(S_6\).
Thus the second evaluation does not reuse the discovery integral.
Last transformed terms are recorded as diagnostics; their small
size alone is not asserted to bound the complete tail.

The weight-seven functional \eqref{eq:separating} sends the
integer-normalized candidate to \(485683200\). Its proof therefore
also requires information beyond the restricted single-product
presentation. A useful next step is to obtain a functional equation
that proves both \eqref{eq:S4-original} and \eqref{eq:S6}, with the
nonzero functional serving as a check that the missing relation has
actually been introduced.

